\documentclass[12pt]{article}
\usepackage[utf8]{inputenc}
\usepackage[T1]{fontenc}
\usepackage{hyperref}
\title{Tikar Language Resources for K-12 Curriculum Development}
\author{Compiled by Claude Sonnet 5 (claude-sonnet-5) for Andr\'as Kornai}
\date{}
\begin{document}
\maketitle

\section{Introduction}
Tikar (ISO 639-3: tik) is a Grassfields Bantu language of west-central Cameroon. Population figures vary widely between sources (about 25,000 for the Bankim community, 50,000 and 184,000 elsewhere), and varieties overlap (Tikar de Bankim, Tikar de Ngambe, Twumwu), so check which variety each resource covers.

\section{Primary Educational Resources}
No primers or school materials were found.

\section{Linguistic and Reference Materials}
Online dictionary of the Bankim variety, with English and French glosses and example sentences, on Webonary: \url{https://www.webonary.work/tikar/overview/introduction/?lang=en}. Foreword: \url{https://www.webonary.org/tikar/overview/foreword/?lang=en}

\section{Religious and Community Resources}
ScriptureEarth lists Tikar scripture (New Testament PDF, ePub, apps, Bible.is links): \url{https://www.scriptureearth.org/00eng.php?iso=tik}. Joshua Project language page: \url{https://www.joshuaproject.net/index.php/languages/tik}; Tikar, Twumwu people group: \url{https://joshuaproject.net/people_groups/15486}

\section{Digital and Audio Resources}
Global Recordings hosts audio Bible stories and messages: \url{https://globalrecordings.net/en/language/tik}. Audio New Testament links are reported via ScriptureEarth.

\section{Curriculum Development Context}
Scripture and a dictionary give a text and vocabulary base for one variety. No school programme or textbooks were found.

\end{document}